% !TeX program = lualatex
% ================================================================
% Bengali Mathematical Dictionary
%
% Generated by the server using the following settings:
%
%   language            Bengali (ben)
%   text direction      left to right
%   font                Noto Sans Bengali
%   fallback font       Noto Serif
%   built from          latex/dictionary/mathematicalDictionary.tex
%   tier 3 vocabulary   green   RGB 0 128 0
%   English text        black   RGB 0 0 0
%   Bengali text        purple  RGB 112 48 160
%   layout macros       version 2
%
% Please don't edit any information in this comment block - the server
% compares it against the current settings and rebuilds this file when
% they differ, so changes here would be overwritten. However, feel free
% to make updates to the LaTeX below if you want to edit it.
% ================================================================
%
% ---- What is safe to edit in this file ----------------------------------
%
% Safe:     the Bengali word and the Bengali definition - the third
%           and fourth arguments of each entry. That is what this file is for:
%           if a translation reads wrongly to somebody who speaks the language,
%           correct it here and every sheet that uses the word follows.
%
% Not safe: the English word and the English definition - the first and second
%           arguments. The first is how an entry is matched to the shared
%           dictionary; the second is the wording the translation was made
%           from, and is how the server knows the translation is still current.
%           Changing an English definition here does not change the shared
%           dictionary - it only makes this entry look up to date when it is
%           not. Edit the English in the shared dictionary instead, and this
%           file will be brought back into step on the next run.
%
% Entries are added and removed by the server, following the shared dictionary.
% An entry the server cannot read is reported rather than guessed at, so if a
% run complains about this file, look for a missing brace.
% -------------------------------------------------------------------------

\documentclass[11pt]{article}
\usepackage[a4paper,margin=18mm]{geometry}
% Characters this language's font has not got are borrowed from another,
% rather than being silently dropped - see docs/translated-sheet-latex.md
\directlua{%
  if luaotfload and luaotfload.add_fallback then
    luaotfload.add_fallback("eallatinfallback",
      { "Noto Serif:mode=harf;" })
  else
    tex.error("this luaotfload is too old to register a fallback font")
  end
}
\usepackage[english]{babel}
\babelprovide[import]{bengali}
\babelfont[bengali]{rm}[Renderer=Harfbuzz, BoldFont={Noto Sans Bengali}, ItalicFont={Noto Sans Bengali}, BoldItalicFont={Noto Sans Bengali}, RawFeature={fallback=eallatinfallback}]{Noto Sans Bengali}
\babelfont[bengali]{sf}[Renderer=Harfbuzz, BoldFont={Noto Sans Bengali}, ItalicFont={Noto Sans Bengali}, BoldItalicFont={Noto Sans Bengali}, RawFeature={fallback=eallatinfallback}]{Noto Sans Bengali}
\newcommand{\ealtext}[1]{\foreignlanguage{bengali}{#1}}
\newcommand{\ealtextblock}[1]{{\raggedright\foreignlanguage{bengali}{#1}\par}}
% ================================================================
% EAL layout helpers
% ================================================================
\usepackage{xcolor}

\definecolor{ealkeycolour}{RGB}{0,128,0}
\definecolor{ealenglish}{RGB}{0,0,0}
\definecolor{ealtwocolour}{RGB}{112,48,160}

\newsavebox{\ealboxen}
\newsavebox{\ealboxtr}
\newlength{\ealglosswd}
\newlength{\ealglossraise}
\setlength{\ealglossraise}{1.15em}

% a tier 3 word, in English and in translation alike
\newcommand{\ealkey}[1]{\textcolor{ealkeycolour}{#1}}
\newcommand{\ealkeytr}[1]{\textcolor{ealkeycolour}{\ealtext{#1}}}

% One translated word above one English word. Built from kernel
% primitives, never a tabular, which would break wherever the sheet
% loads array, booktabs or tabularx. The box is as wide as the wider of
% the two, so a long translation pushes its neighbours aside instead of
% overlapping them, and it claims its full height so a table row or a
% TikZ node makes room for it.
\newcommand{\ealgloss}[2]{%
  \sbox{\ealboxen}{\ealkey{#1}}%
  \sbox{\ealboxtr}{\small\textcolor{ealtwocolour}{\ealtext{#2}}}%
  \setlength{\ealglosswd}{\wd\ealboxen}%
  \ifdim\wd\ealboxtr>\ealglosswd\setlength{\ealglosswd}{\wd\ealboxtr}\fi
  \leavevmode
  \makebox[\ealglosswd][c]{%
    \makebox[0pt][c]{\usebox{\ealboxen}}%
    \makebox[0pt][c]{\raisebox{\ealglossraise}{\usebox{\ealboxtr}}}%
  }%
}

% opens up line spacing around a block of glosses, so the raised
% translations sit clear of the line above
\newenvironment{ealglossed}{\par\linespread{2.1}\selectfont\ignorespaces}{\par}

% a whole translated sentence above its English counterpart. block level,
% so it does not belong inside a tabular cell
\newcommand{\ealpara}[2]{%
  \par\addvspace{0.5\baselineskip}%
  {\color{ealtwocolour}\ealtextblock{#1}}%
  \penalty200
  {\color{ealenglish}#2\par}%
  \addvspace{0.5\baselineskip}%
}

% ================================================================
% Dictionary layout
% ================================================================
\usepackage{multicol}

\newlength{\ealdictgap}
\setlength{\ealdictgap}{1.2mm}

% english word / english meaning / translated word / translated meaning
\newcommand{\dictentrytr}[4]{%
  \par\addvspace{\ealdictgap}%
  \noindent
  {\bfseries\ealkey{#1}}\quad{\small #2}\par
  \nopagebreak
  {\color{ealtwocolour}\ealtext{{\bfseries #3}~\textemdash{} {\small #4}}\par}%
  \par
}

\pagestyle{empty}
\setlength{\parindent}{0pt}

\begin{document}
{\large\bfseries Bengali Mathematical Dictionary}\par\medskip

\dictentrytr{coefficient}{the number in front of a letter, showing how many of it there are}{সহগ}{একটি অক্ষরের সামনের সংখ্যা, যা দেখায় সেটির কতগুলো আছে}
\dictentrytr{decimal}{a number written with a point, where the digits after the point are parts of a whole}{দশমিক}{এমন সংখ্যা যা একটি বিন্দু দিয়ে লেখা হয়, যেখানে বিন্দুর পরে থাকা অঙ্কগুলো একটি পূর্ণের অংশ}
\dictentrytr{denominator}{the number below the line in a fraction, showing how many equal parts the whole is split into}{হর}{ভগ্নাংশে রেখার নিচের সংখ্যা, যা দেখায় পূর্ণটি কতটি সমান ভাগে ভাগ করা হয়েছে}
\dictentrytr{difference}{the answer when one number is subtracted from another}{বিয়োগফল}{একটি সংখ্যা থেকে আরেকটি সংখ্যা বিয়োগ করলে যে উত্তর পাওয়া যায়}
\dictentrytr{digit}{a single figure from 0 to 9 that a number is written with}{অঙ্ক}{০ থেকে ৯ পর্যন্ত একটি একক চিহ্ন, যার সাহায্যে সংখ্যা লেখা হয়}
\dictentrytr{equivalent}{having the same value, even though it is written differently}{সমতুল্য}{একই মান আছে এমন, যদিও এটি ভিন্নভাবে লেখা হয়}
\dictentrytr{estimate}{a sensible answer worked out quickly from rounded numbers}{আনুমানিক মান}{আনুমানিক সংখ্যা থেকে দ্রুত বের করা একটি যুক্তিসঙ্গত উত্তর}
\dictentrytr{expand}{to multiply out brackets so the expression has no brackets left}{প্রসারিত করা}{বন্ধনীগুলো গুণ করে খোলা, যাতে রাশিটিতে আর কোনো বন্ধনী না থাকে}
\dictentrytr{expression}{a group of numbers and letters joined by operations, with no equals sign}{রাশি}{সংখ্যা ও অক্ষরের একটি গুচ্ছ যা গাণিতিক প্রক্রিয়া দিয়ে যুক্ত থাকে এবং যাতে সমান চিহ্ন থাকে না}
\dictentrytr{factor}{a whole number that divides exactly into another number}{গুণনীয়ক}{একটি পূর্ণ সংখ্যা যা অন্য একটি সংখ্যাকে নিঃশেষে ভাগ করে}
\dictentrytr{factorise}{to write an expression as a product, usually by putting brackets back in}{উৎপাদকে বিশ্লেষণ করা}{একটি রাশিকে গুণফল আকারে লেখা, সাধারণত বন্ধনী পুনরায় বসিয়ে}
\dictentrytr{formula}{a rule written with letters, showing how to work one quantity out from others}{সূত্র}{অক্ষর দিয়ে লেখা একটি নিয়ম, যা দেখায় কীভাবে একটি রাশি অন্যগুলো থেকে বের করা যায়}
\dictentrytr{fraction}{a number written as one whole split into equal parts, with a number of those parts taken}{ভগ্নাংশ}{একটি সংখ্যা, যা একটি পূর্ণকে সমান ভাগে ভাগ করে সেগুলোর কয়েকটি নেওয়ার মাধ্যমে লেখা হয়}
\dictentrytr{gradient}{how steep a line is, found by how far it rises for each step across}{ঢাল}{একটি রেখা কতটা খাড়া, তা বের করা হয় রেখাটি প্রতিটি ধাপ পাশে যাওয়ার সময় কতটা উপরে ওঠে তা দেখে}
\dictentrytr{highest common factor}{the largest whole number that divides exactly into two or more numbers}{গরিষ্ঠ সাধারণ গুণনীয়ক}{সবচেয়ে বড় পূর্ণ সংখ্যা যা দুই বা ততোধিক সংখ্যাকে নিঃশেষে ভাগ করে}
\dictentrytr{identity}{a statement that is true for every value of the letters in it}{অভেদ}{এমন একটি বিবৃতি যা এর অক্ষরগুলোর প্রতিটি মানের জন্যই সত্য}
\dictentrytr{index}{the small raised number showing what power something is raised to}{সূচক}{ছোট করে উপরে লেখা সংখ্যা, যা নির্দেশ করে কোনো কিছুকে কী ঘাতে তোলা হয়েছে}
\dictentrytr{inequality}{a statement that one quantity is larger or smaller than another}{অসমতা}{এমন একটি বিবৃতি যা বলে একটি রাশি অন্য রাশির চেয়ে বড় বা ছোট}
\dictentrytr{integer}{a whole number, which may be positive, negative or zero}{পূর্ণসংখ্যা}{একটি পূর্ণ সংখ্যা, যা ধনাত্মক, ঋণাত্মক বা শূন্য হতে পারে}
\dictentrytr{lowest common multiple}{the smallest number that two or more numbers all divide into exactly}{লঘিষ্ঠ সাধারণ গুণিতক}{যে ক্ষুদ্রতম সংখ্যাটিকে দুই বা ততোধিক সংখ্যার প্রত্যেকটি নিঃশেষে ভাগ করে}
\dictentrytr{multiple}{a number you get by multiplying another number by a whole number}{গুণিতক}{অন্য একটি সংখ্যাকে একটি পূর্ণ সংখ্যা দিয়ে গুণ করে যে সংখ্যা পাওয়া যায়}
\dictentrytr{multiply}{to add a number to itself a given number of times}{গুণ করা}{একটি সংখ্যাকে নির্দিষ্ট সংখ্যক বার নিজের সাথে যোগ করা}
\dictentrytr{numerator}{the number above the line in a fraction, showing how many equal parts are taken}{লব}{ভগ্নাংশে রেখার উপরের সংখ্যা, যা দেখায় কতটি সমান ভাগ নেওয়া হয়েছে}
\dictentrytr{percentage}{a number of parts out of a hundred}{শতকরা}{একশোর মধ্যে কতগুলো অংশ}
\dictentrytr{place value}{the value a digit has because of where it sits in a number}{স্থানীয় মান}{একটি অঙ্ক সংখ্যায় যেখানে অবস্থান করে তার জন্য যে মান পায়}
\dictentrytr{power}{how many times a number is multiplied by itself}{ঘাত}{একটি সংখ্যাকে নিজের সাথে কতবার গুণ করা হয়}
\dictentrytr{prime}{a whole number greater than 1 that divides only by itself and 1}{মৌলিক সংখ্যা}{১-এর চেয়ে বড় একটি পূর্ণ সংখ্যা যা শুধু নিজের এবং ১ দ্বারাই বিভাজ্য}
\dictentrytr{product}{the answer when two or more numbers are multiplied}{গুণফল}{দুই বা ততোধিক সংখ্যাকে গুণ করলে যে উত্তর পাওয়া যায়}
\dictentrytr{proportion}{a relationship where two amounts change by the same factor}{সমানুপাত}{এমন সম্পর্ক যেখানে দুটি পরিমাণ একই গুণক দিয়ে পরিবর্তিত হয়}
\dictentrytr{quotient}{the answer when one number is divided by another}{ভাগফল}{একটি সংখ্যাকে আরেকটি সংখ্যা দিয়ে ভাগ করলে যে উত্তর পাওয়া যায়}
\dictentrytr{ratio}{a way of comparing two or more amounts, showing how much of one there is for each of the other}{অনুপাত}{দুই বা ততোধিক পরিমাণ তুলনা করার একটি উপায়, যা দেখায় অন্যগুলোর প্রতিটির জন্য একটির কতটুকু আছে}
\dictentrytr{reciprocal}{the number you multiply by to get 1, found by turning a fraction upside down}{ব্যাস্ত}{যে সংখ্যা দিয়ে গুণ করলে ১ পাওয়া যায়, এটি ভগ্নাংশটিকে উল্টে দিলে পাওয়া যায়}
\dictentrytr{recur}{to repeat the same digits forever after the decimal point}{পুনরাবৃত্ত হওয়া}{দশমিক বিন্দুর পরে একই অঙ্কগুলো চিরকাল পুনরাবৃত্ত হতে থাকে}
\dictentrytr{root}{a number that gives another number when multiplied by itself a given number of times}{মূল}{এমন একটি সংখ্যা যা নির্দিষ্ট সংখ্যক বার নিজের সাথে গুণ করলে আরেকটি সংখ্যা পাওয়া যায়}
\dictentrytr{round}{to replace a number with a nearby simpler one}{রাউন্ড করা}{একটি সংখ্যাকে তার নিকটবর্তী সরল সংখ্যা দিয়ে প্রতিস্থাপন করা}
\dictentrytr{scale factor}{the number every length is multiplied by when a shape is enlarged}{মাপনী গুণক}{একটি আকার বড় করার সময় প্রতিটি দৈর্ঘ্যকে যে সংখ্যা দিয়ে গুণ করা হয়}
\dictentrytr{simplify}{to write something in its smallest or plainest form without changing its value}{সরল করা}{মান না বদলে কোনো কিছুকে তার ক্ষুদ্রতম বা সরলতম আকারে লেখা}
\dictentrytr{sum}{the answer when numbers are added together}{যোগফল}{সংখ্যাগুলো একসাথে যোগ করলে যে উত্তর পাওয়া যায়}
\dictentrytr{surd}{a root that cannot be written exactly as a fraction, so it is left in root form}{করণী}{এমন একটি মূল যা ভগ্নাংশ হিসেবে সঠিকভাবে লেখা যায় না, তাই একে মূল আকারে রাখা হয়}
\dictentrytr{unitary}{a method that first finds the value of one part or one item}{একক পদ্ধতি}{এমন একটি পদ্ধতি যেখানে প্রথমে একটি অংশ বা একটি বস্তুর মান বের করা হয়}

\end{document}
